\documentclass[12pt]{book}
\usepackage{amsmath}
\usepackage{amssymb}
\usepackage{geometry}
\geometry{margin=1in}

\title{Volume 07: Gravitation and Stellar Structure \\ Book 01: Gravitation from Matrix Pressure}
\author{James C. Harvey}
\date{December 2025}

\begin{document}
\maketitle
\tableofcontents

\chapter{Gravity as a Regime}

\section*{Abstract}
Gravitation is the large-boundary, low-curvature regime of the universal coupling equation. It is not a
separate force but a scale limit of flux geometry.

\section{Regime Expression}
For large boundaries,
\begin{equation}
F_{\text{grav}} = \frac{\kappa_1 \kappa_2}{4\pi K_{\text{bulk}} r^2}\,\mathcal{C}(\mathbf{r},\hat{\mathbf{n}}).
\end{equation}
The composite constant is
\begin{equation}
G = \frac{c^4}{4\pi K_{\text{bulk}}^2}.
\end{equation}

\chapter{Stellar Structure from Pressure Geometry}

\section*{Abstract}
Stars are equilibrium structures where flux pressure gradients balance boundary curvature and internal
motion. This chapter encodes the geometric balance conditions.

\section{Balance Equation}
\begin{equation}
\frac{\kappa_{\text{star}} \kappa_{\text{shell}}}{r^2} = \frac{\kappa_{\text{shell}} v^2}{r}.
\end{equation}

\chapter{Gravitational Potential}

\section*{Abstract}
Gravitational potential is the large-scale pressure field. It is a map of flux deficit created by
macroscopic boundary geometry.

\section{Potential Definition}
\begin{equation}
\Phi_G(\mathbf{r}) = -\int_{\infty}^{\mathbf{r}} \frac{\kappa_{\text{eff}}}{4\pi K_{\text{bulk}} r'^2}\, dr'.
\end{equation}

\chapter{Orbital Dynamics}

\section*{Abstract}
Orbits are stable paths in a pressure gradient. This chapter expresses orbital speed and stability in
SDT terms.

\section{Circular Orbit}
\begin{equation}
v^2 = \frac{\kappa_1 \kappa_2}{4\pi K_{\text{bulk}} r}.
\end{equation}

\chapter{Gravitational Waves as Flux Oscillation}

\section*{Abstract}
Gravitational waves are oscillatory pressure modes in the matrix. They propagate as large-scale flux
oscillations at the matrix wave speed.

\section{Wave Equation}
\begin{equation}
\frac{\partial^2 \Pi}{\partial t^2} = c^2 \nabla^2 \Pi.
\end{equation}

\chapter{Stellar Evolution Regimes}

\section*{Abstract}
Evolution is the change of boundary curvature and flux throughput. This chapter summarizes the SDT
view of stellar stages as geometry transitions.

\section{Geometry Transition}
\begin{equation}
\Delta \kappa \propto \Delta A_{\text{eff}}^{-1},
\end{equation}
so that contraction increases curvature and alters flux capture.
\end{document}
